\section{Overview and Key Ideas}
\label{sec:overview-and-key-ideas}

To address the two key inefficiencies identified in Section~\ref{sec:motivation}: runtime overhead from redundant memory access and throughput degradation due to load imbalance and parallelism overhead, we propose \textit{\Name}, a novel LLM LoRA fine-tuning system that improves system throughput via kernel-level optimizations and job-level scheduling.
In this section, we summarize the two key ideas behind it:

\begin{figure}[t]
    \centering
    \includegraphics[width=0.975\linewidth]{figures/overview.pdf}
    \caption{Overview of LoRAFusion.}
    \label{fig:key-ideas-overview}
\end{figure}

\myparagraph{1. FusedLoRA and FusedMultiLoRA}
To reduce LoRA's runtime overhead, \Name fuses memory-bound operations while preserving performance for compute-bound matrix multiplications by carefully splitting the computation graph.
This is motivated by the key bottleneck in LoRA, i.e., redundant memory accesses to large activation tensors. 
A naive solution is to fuse the entire computation graph into a single kernel. 
However, doing so introduces costly recomputation or synchronization overhead~\cite{shi2023welder}. It also consumes GPU resources like registers and shared memory, which degrades GEMM performance due to suboptimal tiling.

\Name addresses this by introducing a graph-splitting strategy. 
Instead of full fusion, it splits the graph at the intermediate tensors with LoRA rank $r$, which are small and cheap to materialize. 
This enables fusion around full-sized activations without recomputation or synchronization.
Thus, the resulting FusedLoRA kernels reduce memory traffic while preserving the performance of frozen GEMMs.

We further extend this idea to support multiple LoRA adapters using a tile-level routing mechanism.
The proposed FusedMultiLoRA processes tokens from different jobs within the same fused kernel.
Each tile uses a precomputed mapping to dynamically select the appropriate adapter weights.
This avoids the need for separate kernel launches per adapter and maintains high GPU utilization across jobs.
During backpropagation, gradients are routed similarly. The FusedMultiLoRA forms the foundation of \Name's job-level scheduling.



\myparagraph{2. Multi-LoRA Scheduling}
To exploit the benefits of multi-LoRA fine-tuning, \Name introduces a scheduler that coordinates adapter grouping and data batching.
The major challenge is to construct well-balanced microbatches across jobs while respecting data dependencies between global batches. Specifically, in pipeline parallelism, each global batch must wait for all samples from the previous one to complete their backward passes.
\Name addresses this with a two-stage hierarchical strategy: it first groups adapters in a way that keeps their batches spread apart in the schedule, then solves a bin-packing problem to batch their samples and reduce imbalance.

In the first stage, \Name groups LoRA adapters based on their sample length distributions to ensure that consecutive global batches from the same adapter are sufficiently spaced in the schedule.
In the second stage, samples within each group are packed into microbatches using a two-step MILP-based optimization, with the greedy algorithm as a fallback and multiprocessing for efficiency.
A final merge pass reduces underfilled microbatches when possible.
This scheduling strategy improves load balance and reduces pipeline stalls, boosting overall system throughput.




% To exploit multi-LoRA fine-tuning, \Name introduces a scheduler that coordinates adapter grouping and sample batching. 
% Each fine-tuning job comes with its own dataset, often with distinct token-length distributions. 
% To minimize token imbalance across microbatches and improve load balance and pipeline efficiency, \Name first groups adapters with complementary sample lengths. 
% This grouping follows a "head-tail" pairing strategy based on statistical properties (e.g., mean and variance) of each dataset, ensuring more uniform microbatches when grouped.
% Within each group, \Name batches samples using a two-stage MILP-based strategy. 
% For each aggregated global batch from grouped adapters, \Name first determines the number of microbatches under a fixed token budget, then assigns samples to minimize token imbalance. 
% If the MILP solver times out, a greedy fallback is used. Groups are scheduled in parallel using multiprocessing, and a post-processing pass merges underfilled microbatches if doing so satisfies pipeline constraints. These strategies jointly reduce pipeline bubbles and improve device utilization.

\myparagraph{System Workflow}
Figure~\ref{fig:key-ideas-overview} presents \Name's system workflow. Given a set of fine-tuning jobs, \Name first extracts dataset statistics and proposes a microbatch token budget via a parallelism simulator. It then forms adapter groups and constructs microbatches accordingly. The grouping and batching outputs are re-evaluated through simulation, and the process iterates until a high-throughput configuration is found. Finally, jobs are executed using the fused kernels described above. A multi-adapter runtime coordinator ensures token-to-adapter consistency, manages resource sharing, and tracks gradients across job boundaries.
Through the combination of fused execution and coordinated job scheduling, \Name addresses both the memory and parallelism bottlenecks in LoRA fine-tuning, improving throughput while maintaining correctness and generality.

